\chapter{Ecosystem Modeling Using Cellular Automata}
\label{ch:ecosystem}

\section{CAecosystem}

The module \texttt{cadyn.ecosystem} illustrates how CA can create
simple, intuitive models of complex real-world phenomena that help
illuminate the underlying principles.  This ecosystem simulator lets
the user set the parameters of an artificial ecosystem with three
interdependent species --- \emph{plants}, \emph{prey}, and
\emph{predators} --- then set the artificial universe in motion,
observe the ensuing dynamics, and display the population data in
various ways.  These created worlds have a toroidal topology: a critter
creeping off the bottom of the display reappears at the top, and one
crawling off the right enters on the left.  So don't worry --- these
critters can't escape and infest your home!

For a first run we use a $500\times500$ cell space with the initial
probability of a cell being prey set to $0.212$ and of being a predator
set to $0.029$; other parameters keep their default settings.
Figure~\ref{fig:eco-init} shows the initial state
(green $=$ plants, yellow $=$ prey, red $=$ predators), and
Figure~\ref{fig:eco-2000} shows the state after 2000 generations
(purple cells are devoid of life).

\begin{figure}[htbp]\centering
\includegraphics[width=0.7\textwidth]{ch5_eco_initial}
\caption{Initial state of the ecosystem (detail).  Green $=$ plants,
yellow $=$ prey, red $=$ predators.}
\label{fig:eco-init}
\end{figure}

\begin{figure}[htbp]\centering
\includegraphics[width=0.7\textwidth]{ch5_eco_gen2000}
\caption{The same ecosystem 2000 generations later (detail); purple
cells are devoid of life.}
\label{fig:eco-2000}
\end{figure}

Figure~\ref{fig:eco-pops} graphs the populations against time, up to
generation 2000.  Without the regulating influence of the predators,
the prey would eat all the plants and then starve.  The three
populations are graphed on a normalized scale for good resolution on a
shared axis.  The predator curve (red) and the prey curve (yellow) show
the out-of-phase oscillation typical of predator/prey systems: abundant
prey mean easy hunting, hence many predator births; the swelling
predator population then drives the prey into decline; scarce prey then
starve the predators; and with predators scarce, the prey rebound,
restarting the cycle.

\begin{figure}[htbp]\centering
\includegraphics[width=0.9\textwidth]{ch5_eco_populations}
\caption{Normalized population curves; the horizontal axis is the
generation number.}
\label{fig:eco-pops}
\end{figure}

The cycle is less obvious between prey and plants (relative to which the
prey are the predators), because of the rapid growth of plants in this
simulation; the plant graph is almost the negative of the prey graph
(roughly $180^\circ$ out of phase).  Notice that the initially large
oscillations (\emph{transients}) settle to relatively stable
amplitudes for all species, and that the predators and plants stand in
a curiously symbiotic relationship in this ecosystem.

Another view of the same information is the 3-dimensional \emph{phase
portrait}: at each sampled time we plot a point whose coordinates are
the three population counts (Figure~\ref{fig:eco-phase}).  Poincar\'e
introduced this type of graph in his study of the global dynamics of
differential equations.  The oscillation of Figure~\ref{fig:eco-pops}
appears here as an almost circular orbiting of the curve about an
equilibrium point, and the damping of the oscillations appears as a
tightening of the spiral toward that point.

\begin{figure}[htbp]\centering
\includegraphics[width=0.7\textwidth]{ch5_eco_phase}
\caption{A 3-dimensional phase portrait representing the same data as
Figure~\ref{fig:eco-pops}: plants ($x$), prey ($y$), predators
($z$).}
\label{fig:eco-phase}
\end{figure}

This same qualitative behavior appears in solutions of the
Volterra--Lotka differential equations traditionally used to model
ecosystem population dynamics.  A little experimentation reveals that an
imprudent choice of parameters produces transients so wild that all
animal life (often all plant life too) goes extinct.
Figure~\ref{fig:extinction} shows such a run: with voracious,
short-lived predators (high predation efficiency, high death rate), the
first predator boom overshoots so violently that it eats the prey
faster than the prey can be replenished, and having exhausted its food
supply the predator population crashes to zero.  Once an animal species
is gone it cannot return --- there is no spontaneous generation in this
model --- so the extinction is permanent.  Even in the
differential-equation models the equilibrium point is virtually
impossible to attain, since it would almost always entail fractions of
an animal being alive --- one of the absurdities of using continuous
models for discrete phenomena.  The difficulty is even more acute in
discrete models, so oscillatory behavior is unavoidable.

\begin{figure}[htbp]\centering
\includegraphics[width=0.85\textwidth]{ch5_extinction}
\caption{An imprudent parameter choice: with predators too efficient
and too short-lived, the first boom-and-bust transient drives the
predators (and here the prey) extinct within the first hundred
generations.}
\label{fig:extinction}
\end{figure}

Which parameter regimes permit long-term coexistence, and which lead to
collapse?  Because the model is stochastic and cheap to run, we can map
this out directly rather than analyze it.  Figure~\ref{fig:survival}
sweeps two predator parameters --- predation efficiency
$b_{\mathrm{pred}}$ and death rate $d_{\mathrm{pred}}$ --- and records
the fraction of runs in which all three species survive 400
generations.  A broad green region of coexistence separates two modes
of failure: predators too feeble or too mortal (lower right) starve and
die out, while predators too voracious and too persistent (upper left)
strip the ecosystem in a single overshoot.  Persistent predator/prey
coexistence is not the generic outcome but a balance struck within a
bounded region of parameter space --- the discrete, spatial echo of the
delicate neutral cycles of the Volterra--Lotka equations.

\begin{figure}[htbp]\centering
\includegraphics[width=0.75\textwidth]{ch5_survival_sweep}
\caption{Survival phase diagram: the fraction of runs (of a small
ensemble) in which all species persist for 400 generations, as the
predator's predation efficiency and death rate are varied.  Green marks
coexistence; red marks extinction.}
\label{fig:survival}
\end{figure}

One reason for the initial transients is that the initial spatial
arrangement of cell states is random, whereas the steady-state
distribution is far from random; the transient period involves the
spatial drift from the random initial configuration toward a more
optimal one (compare Figures~\ref{fig:eco-init}
and~\ref{fig:eco-2000}).  Another reason, visible in the phase
portrait, is that our initial state was a bit too lush with plant life,
placing the initial point off to one side of the main coil of the
trajectory.

\expand{This CA ecosystem is a spatial, individual-based analogue of
the Volterra--Lotka equations, and it illustrates a lesson that the
intervening decades have underscored: spatial structure changes
everything.  Where the mean-field differential equations predict either
neutral cycles or a single damped approach to equilibrium, spatially
extended predator/prey CA generate travelling waves, spiral patterns,
and persistent noisy oscillations sustained by local
desynchronization --- dynamics now central to theoretical ecology and
epidemiology.  The same mechanism drives spiral waves in the
Belousov--Zhabotinsky reaction and in models of cardiac fibrillation
mentioned in Chapter~1.  This ecosystem is also the springboard for
Chapter~\ref{ch:evolution}, which gives its creatures the ability to
\emph{move} and equips them with competing perceptual \emph{strategies},
turning the simulation into an evolutionary-game laboratory.}

The original edition tracked the same ecosystem out to 11{,}200
generations; the population curves continue their noisy, bounded
oscillation indefinitely, and the actual state continues to resemble
Figure~\ref{fig:eco-2000}.  The reader is
encouraged to reproduce such long runs --- now a matter of seconds
rather than an overnight computation --- and to rotate the phase
portrait interactively:

\begin{lstlisting}
from cadyn.ecosystem import Ecosystem
import numpy as np

eco = Ecosystem(n=200, p_prey=0.212, p_pred=0.029,
                rng=np.random.default_rng(11))
eco.run(11_200)
eco.plot_populations(sample_every=8)   # long-run population curves
ax = eco.plot_phase_portrait()         # rotate with ax.view_init(...)
\end{lstlisting}
